% -----------------------------------------------------------------
% Compact L&FW section, mirroring Mai & Lin (TMLR 2026, §6) — two
% concrete limitations with future-work hooks woven in. ICML hooks
% kept: univariate scope, single global α-factor.
% -----------------------------------------------------------------
\section{Limitations and Future Work}
\label{sec:limitations}

% REVISED FOR TMLR (rebuttal C2)
Despite its contributions, this study has three limitations. First,
our evaluation is restricted to \emph{univariate} TS;
extending the contrastive clustering loss to a joint-channel latent
space, where inter-channel dependencies are modelled directly, is
the natural next step toward multivariate sensor data. Second,
the expansion uses a single scalar $\alpha$ applied uniformly across all
classes. Alternative expansion mechanisms that adapt to class-specific
covariance structure would relax this restriction and may further
improve gains on heterogeneous dataset families.
% ADDED FOR TMLR (rebuttal C2)
Third, the benefit is uneven across classifiers. FCN is the
weakest case, with the median dataset unchanged in all three
fixed-configuration runs, so on FCN we describe ASCENSION as
low-risk rather than as broadly beneficial.

% ADDED FOR TMLR (rebuttal Q5/Q6)
Beyond these limitations, two extensions follow naturally from the
design. ASCENSION could handle regression targets by conditioning
the latent space on the continuous target, or by binning the
target into pseudo-classes that take the role of labels in the
contrastive clustering loss. It could also apply beyond univariate
TS to other input domains, since only the encoder is
domain-specific. We have evaluated neither extension and claim no
benefit for either without supporting evidence.
